% Part: first-order-logic
% Chapter: syntax-and-semantics
% Section: substitution

\documentclass[../../../include/open-logic-section]{subfiles}

\begin{document}

\olfileid{fol}{syn}{sub}

\olsection{تھاں بدل}

\begin{defn}[ٹرم وچ تھانواں بدلنا]
اسیں $\Subst{s}{t}{x}$ نوں، یعنی $t$ نوں~$x$ دی
ہر آمد دی تھاں~$s$ وچ \emph{رکھن} دا نتیجہ، بازگشتی
طور تے تعریف کردے آں:
\begin{enumerate}
\item \indcase{s}{c}{$\Subst{\indfrm}{t}{x}$ بس $s$ اے۔}

\item \indcase{s}{y}{$\Subst{\indfrm}{t}{x}$ وی بس~$s$ اے،
  بشرطیکہ $y$ متغیر ہووے تے $y \not\ident x$ ہووے۔}

\item \indcase{s}{x}{$\Subst{\indfrm}{t}{x}$،~$t$ اے۔}

\item\indcase{s}{\Atom{f}{t_1, \dots, t_n}}{$\Subst{\indfrmp}{t}{x}$،
  $\Atom{f}{\Subst{t_1}{t}{x}, \dots, \Subst{t_n}{t}{x}}$ اے۔}
\end{enumerate}
\end{defn}

\begin{defn}
ٹرم~$t$، $x$ لئی~$!A$ وچ \emph{!!{free for}} اے
جے~$x$ دی کوئی آزاد آمد~$!A$ وچ کسے ایسے مسوّر دے
دائرۂ اثر وچ نہ آوے جیہڑا~$t$ دے کسے متغیر نوں پابند کردا ہووے۔
\end{defn}

\begin{ex} ~
\begin{enumerate}
\item $\Obj v_8$، $\Obj v_1$ لئی
$\lexists[\Obj
  v_3]\Atom{\Obj A^2_4}{\Obj v_3,\Obj v_1}$ وچ بدل لئی آزاد اے۔

\item $\Obj f^2_1(\Obj v_1, \Obj v_2)$، $\Obj
  v_0$ لئی $\lforall[\Obj v_2]\Atom{\Obj A^2_4}{\Obj v_0,\Obj v_2}$
  وچ بدل لئی \emph{آزاد نہیں} اے۔
\end{enumerate}
\end{ex}

\begin{defn}[!!a{formula} وچ تھانواں بدلنا]
جے $!A$، !!a{formula} اے، $x$~!!a{variable} اے،
تے $t$~اوہ ٹرم اے جیہڑا~$x$ لئی~$!A$ وچ
!!{free for} اے، تاں $\Subst{!A}{t}{x}$،
$t$ نوں~$x$ دیاں ساریاں آزاد آمدواں دی تھاں
$!A$ وچ رکھن دا نتیجہ اے۔
\begin{enumerate}
\tagitem{prvFalse}{\indcase{!A}{\lfalse}{$\Subst{\indfrm}{t}{x}$،
    $\lfalse$ اے۔}}{}

\tagitem{prvTrue}{\indcase{!A}{\ltrue}{$\Subst{\indfrm}{t}{x}$،
    $\ltrue$ اے۔}}{}

\item \indcase{!A}{\Atom{P}{t_1,\dots,
    t_n}}{$\Subst{\indfrm}{t}{x}$،
    $\Atom{P}{\Subst{t_1}{t}{x},
    \dots, \Subst{t_n}{t}{x}}$ اے۔}

\item \indcase{!A}{\eq[t_1][t_2]}{$\Subst{\indfrmp}{t}{x}$،
  $\Subst{t_1}{t}{x} = \Subst{t_2}{t}{x}$ اے۔}

\tagitem{prvNot}{\indcase{!A}{\lnot !B}{$\Subst{\indfrmp}{t}{x}$،
    $\lnot \Subst{!B}{t}{x}$ اے۔}}{}

\tagitem{prvAnd}{\indcase{!A}{(!B \land
    !C)}{$\Subst{\indfrmp}{t}{x}$،
    $(\Subst{!B}{t}{x} \land
    \Subst{!C}{t}{x})$ اے۔}}{}

\tagitem{prvOr}{\indcase{!A}{(!B \lor
    !C)}{$\Subst{\indfrmp}{t}{x}$،
    $(\Subst{!B}{t}{x} \lor
    \Subst{!C}{t}{x})$ اے۔}}{}

\tagitem{prvIf}{\indcase{!A}{(!B \lif
    !C)}{$\Subst{\indfrmp}{t}{x}$،
    $(\Subst{!B}{t}{x} \lif
    \Subst{!C}{t}{x})$ اے۔}}{}

\tagitem{prvIff}{\indcase{!A}{(!B \liff
    !C)}{$\Subst{\indfrmp}{t}{x}$،
    $(\Subst{!B}{t}{x} \liff
    \Subst{!C}{t}{x})$ اے۔}}{}

\tagitem{prvAll}{
\indcase{!A}{\lforall[y][!B]}{$\Subst{\indfrmp}{t}{x}$
    دا نتیجہ $\lforall[y][\Subst{!B}{t}{x}]$ اے،
    بشرطیکہ $y$، $x$ توں وکھرا متغیر ہووے؛
    نہیں تاں $\Subst{\indfrmp}{t}{x}$
    بس $\indfrm$ اے۔}}{}

\tagitem{prvEx}{
\indcase{!A}{\lexists[y][!B]}{$\Subst{\indfrmp}{t}{x}$
    دا نتیجہ $\lexists[y][\Subst{!B}{t}{x}]$ اے،
    بشرطیکہ $y$، $x$ توں وکھرا متغیر ہووے؛
    نہیں تاں $\Subst{\indfrmp}{t}{x}$
    بس $\indfrm$ اے۔}}{}
\end{enumerate}
\end{defn}

\begin{explain}
دھیان دیو کہ تھانواں بدلنا بے اثر وی ہو سکدا اے:
جے $x$، $!A$ وچ کِتے وی نہ آوے، تاں
$\Subst{!A}{t}{x}$ بس~$!A$ اے۔

$t$ دا~$x$ لئی~$!A$ وچ !!{free for} ہونا
اگلی مثال ورگیاں صورتاں کڈھن لئی ضروری اے۔
جے $!A \ident \lexists[y][x < y]$
تے $t \ident y$، تاں $\Subst{!A}{t}{x}$
دا نتیجہ $\lexists[y][y <
  y]$ بن جاندا۔ ایتھے آزاد متغیر $y$،
تھانواں بدلّن توں بعد مسوّر $\lexists[y]$ دی
پابندی وچ ”پھس“ جاندا اے، جیہڑی گل ساڈے لئی
ناپسندیدہ اے۔ مثال لئی، اسیں چاہندے آں کہ
جدوں وی $\lforall[x][!B]$ پورا اترے، تاں
$\Subst{!B}{t}{x}$ وی پورا اترے۔ پر
$\lforall[x][\lexists[y][x < y]]$ ویکھو
(ایتھے $!B$، $\lexists[y][x <
  y]$ اے)۔ ایہ ایسا جملہ اے جیہڑا، مثال لئی،
قدرتی عدداں بارے سچ اے: ہر عدد~$x$ لئی اوہدے
توں وڈّا کوئی عدد~$y$ موجود اے۔ جے اسیں
$y$ نوں~$x$ دی تھاں رکھن دی اجازت دیندے،
تاں $\Subst{!B}{y}{x} \ident \lexists[y][y < y]$
ملدا، جیہڑا جھوٹ اے۔ اسیں ایہ صورت
ایہ شرط رکھ کے روکدے آں کہ~$t$ وچ موجود
کوئی وی آزاد متغیر~$!A$ دے کسے مسوّر نال
پابند نہ ہو جائے۔
\end{explain}

بوجھل علامت نگاری توں بچن لئی اسیں اکثر ایہ رواج
ورتدے آں: جے $!A$ اک !!a{formula} اے جیہدے وچ
!!{variable}~$x$ آزاد ہو سکدا اے، تاں
اسیں ایہ وکھاؤن لئی~$!A(x)$ وی لکھدے آں۔
جدوں ایہ صاف ہووے کہ کیہڑا $!A$ تے~$x$
مراد نیں، تے $t$ اک ٹرم ہووے (فرض کیتا
جاندا اے کہ اوہ $x$ لئی~$!A(x)$ وچ بدل لئی
آزاد اے)، تاں اسیں $!A(t)$ نوں
$\Subst{!A}{t}{x}$ دا مختصر روپ لکھدے آں۔
مثال لئی، اسیں آکھ سکدے آں، ”$!A(t)$ نوں
$\lforall[x][!A(x)]$ دی اک مخصوص صورت آکھدے آں۔“
ایس دا مطلب اے کہ جے $!A$ کوئی وی
!!{formula} ہووے، $x$~!!a{variable} ہووے،
تے $t$ اک ایسا ٹرم ہووے جیہڑا~$x$ لئی~$!A$ وچ
بدل لئی آزاد ہووے، تاں $\Subst{!A}{t}{x}$
فارمولا~$\lforall[x][!A]$ دی اک مخصوص صورت اے۔

\end{document}
